\section{SAC：Soft Actor-Critic}

\subsection{更彻底的 Off-Policy}

PPO 仍然需要定期重新收集数据。SAC 走得更远：使用 \textbf{replay buffer} 存储所有历史数据，从中随机采样来更新 Q 函数和策略。

\subsection{拟合 Q 函数}

SAC 直接拟合 $Q^\pi$ 而不是 $V^\pi$，使用 TD 更新：
$$\min_\phi \sum_{(s,a,s') \in \mathcal{D}} \left\| \hat{Q}^\pi_\phi(s, a) - \left(r(s,a) + \gamma \mathbb{E}_{a' \sim \pi_\theta(\cdot|s')}\left[\hat{Q}^\pi_{\phi'}(s', a')\right]\right)\right\|^2$$

\begin{knowledgebox}{为什么 Q 函数可以用 off-policy 数据}
Q 函数的 TD target $r(s,a) + \gamma \mathbb{E}_{a' \sim \pi}\left[Q(s', a')\right]$ 中，$a'$ 是从\textbf{当前}策略采样的（而不是从数据中读取的）。因此 $(s, a, s')$ 转移数据可以来自任何策略，只要 $a'$ 是新采样的即可。
\end{knowledgebox}

\subsection{策略更新}

策略优化目标加入了\textbf{熵正则化}：
$$\max_\theta \mathbb{E}_{s \sim \mathcal{D}, a \sim \pi_\theta(\cdot|s)} \left[\hat{Q}^\pi(s, a) - \alpha \log \pi_\theta(a \mid s)\right]$$

熵项 $-\alpha \log \pi_\theta$ 鼓励策略保持一定的随机性，有助于探索。

\subsection{SAC 完整算法}

\begin{enumerate}
    \item 用当前策略与环境交互，将转移 $(s, a, r, s')$ 存入 replay buffer $\mathcal{B}$
    \item 从 $\mathcal{B}$ 采样 mini-batch
    \item 更新 Q 函数（用 TD target）
    \item 更新策略（最大化 Q + 熵）
    \item 重复步骤 1-4
\end{enumerate}

\begin{importantbox}{PPO vs. SAC}
\begin{itemize}
    \item \textbf{PPO}：On-policy（加少量 off-policy），用 clipped surrogate objective。每次收集新数据后做几步更新。更稳定、更易调参。
    \item \textbf{SAC}：完全 off-policy，用 replay buffer + Q 函数。数据效率更高。适合连续动作空间。
\end{itemize}
两者都是目前最常用的深度 RL 算法。
\end{importantbox}

\subsection{本章小结}

SAC 通过拟合 Q 函数并使用 replay buffer 实现了完全 off-policy 的 actor-critic 算法，数据效率远高于 on-policy 方法。熵正则化鼓励探索并提高稳定性。

